% ----------------------------------------------------------------- %
%             The Speech Signal Processing Toolkit (SPTK)           %
%             developed by SPTK Working Group                       %
%             http://sp-tk.sourceforge.net/                         %
% ----------------------------------------------------------------- %
%                                                                   %
%  Copyright (c) 1984-2007  Tokyo Institute of Technology           %
%                           Interdisciplinary Graduate School of    %
%                           Science and Engineering                 %
%                                                                   %
%                1996-2016  Nagoya Institute of Technology          %
%                           Department of Computer Science          %
%                                                                   %
% All rights reserved.                                              %
%                                                                   %
% Redistribution and use in source and binary forms, with or        %
% without modification, are permitted provided that the following   %
% conditions are met:                                               %
%                                                                   %
% - Redistributions of source code must retain the above copyright  %
%   notice, this list of conditions and the following disclaimer.   %
% - Redistributions in binary form must reproduce the above         %
%   copyright notice, this list of conditions and the following     %
%   disclaimer in the documentation and/or other materials provided %
%   with the distribution.                                          %
% - Neither the name of the SPTK working group nor the names of its %
%   contributors may be used to endorse or promote products derived %
%   from this software without specific prior written permission.   %
%                                                                   %
% THIS SOFTWARE IS PROVIDED BY THE COPYRIGHT HOLDERS AND            %
% CONTRIBUTORS "AS IS" AND ANY EXPRESS OR IMPLIED WARRANTIES,       %
% INCLUDING, BUT NOT LIMITED TO, THE IMPLIED WARRANTIES OF          %
% MERCHANTABILITY AND FITNESS FOR A PARTICULAR PURPOSE ARE          %
% DISCLAIMED. IN NO EVENT SHALL THE COPYRIGHT OWNER OR CONTRIBUTORS %
% BE LIABLE FOR ANY DIRECT, INDIRECT, INCIDENTAL, SPECIAL,          %
% EXEMPLARY, OR CONSEQUENTIAL DAMAGES (INCLUDING, BUT NOT LIMITED   %
% TO, PROCUREMENT OF SUBSTITUTE GOODS OR SERVICES; LOSS OF USE,     %
% DATA, OR PROFITS; OR BUSINESS INTERRUPTION) HOWEVER CAUSED AND ON %
% ANY THEORY OF LIABILITY, WHETHER IN CONTRACT, STRICT LIABILITY,   %
% OR TORT (INCLUDING NEGLIGENCE OR OTHERWISE) ARISING IN ANY WAY    %
% OUT OF THE USE OF THIS SOFTWARE, EVEN IF ADVISED OF THE           %
% POSSIBILITY OF SUCH DAMAGE.                                       %
% ----------------------------------------------------------------- %
\hypertarget{grlogsp}{}
\name{grlogsp}{draw a running log spectrum graph}{plotting graphs}

\begin{synopsis}
\item[grlogsp] [ --t ] [ --F $F$] [ --O $O$ ] [ --x $X$ ] [ --y $ymin$ ] [ --yy $YY$ ]
	       [ --yo $YO$ ] [ --p $P$ ] 
\item[\ ~~~~~~~~] [ --ln $LN$ ] [ --s $S$ ] [ --e $E$ ] [ --n $N$ ] [ --l $L$ ] 
\item[\ ~~~~~~~~] [ --c $comment1$ ] [ --c2 $comment2$ ] [ --c3 $comment3$ ]
		  [ {\em infile} ]
\end{synopsis}

\begin{qsection}{DESCRIPTION}
{\em grlogsp} converts a sequence of float-format log spectra 
from {\em infile} (or standard input) 
to a running spectrum plot in FP5301 plot format,
sending the result to standard output. 
The output can be visualized with \hyperlink{xgr}{xgr}.

{\em grlogsp} is implemented as a shell script 
that uses the \hyperlink{fig}{fig} and \hyperlink{fdrw}{fdrw} commands.
\end{qsection}

\begin{options}
	\argm{t}{}{transpose x and y axes}{FALSE}
        \argm{F}{F}{factor}{1}
	\argm{O}{O}{origin of graph\\
                    if $O$ is more than 6, drawing area is over A4 range\\
		      \begin{minipage}{4cm}
		      \begin{tabular}{ccc}
			1 & ( 25,$YO$) & [mm] \\
			2 & ( 60,$YO$) & [mm] \\
			3 & ( 95,$YO$) & [mm] \\
			4 & (130,$YO$) & [mm] \\
			5 & (165,$YO$) & [mm] \\
			6 & (200,$YO$) & [mm] \\
			7 & (235,$YO$) & [mm] \\
			8 & (270,$YO$) & [mm] \\
			9 & (305,$YO$) & [mm] \\
		       10 & (340,$YO$) & [mm]
		      \end{tabular}\\\hspace*{\fill}
		      \end{minipage}
		      \begin{minipage}{4cm}
			    \includegraphics{fig/grlogsp-on.pdf}
		      \end{minipage}\\\hspace*{\fill}\\
		    $(YO + 100 , X)$ [mm] if -t is specified.}{1}
	\argm{x}{X}{ $x$ scale\\
		       \begin{tabular}{cl}
			1 & normalized frequency ($0 \sim 0.5$) \\
			2 & normalized frequency ($0 \sim \pi$) \\
			4 & frequency ($0 \sim 4$ kHz) \\
			5 & frequency ($0 \sim 5$ kHz) \\
			8 & frequency ($0 \sim 8$ kHz) \\
			10 & frequency ($0 \sim 10$ kHz) \\
                        12 & frequency ($0 \sim 12$ kHz) \\
			16 & frequency ($0 \sim 16$ kHz) \\
			22 & frequency ($0 \sim 22$ kHz) \\
			24 & frequency ($0 \sim 24$ kHz) \\
			48 & frequency ($0 \sim 48$ kHz) \\
                        96 & frequency ($0 \sim 96$ kHz) \\
                        192 & frequency ($0 \sim 192$ kHz) \\
		       \end{tabular}\\\hspace*{\fill}}{1}
	\argm{y}{ymin}{ $y$ minimum}{-100}
	\argm{yy}{YY}{ $y$ scale [dB/10mm]}{100}
	\argm{yo}{YO}{ $y$ offset}{30}
	\argm{p}{p}{type of pen ($1 \sim 10$)}{2}
	\argm{ln}{LN}{style of line ($0 \sim 5$) (see also \hyperlink{fig}{fig})}{1}
	\argm{s}{S}{start frame number}{0}
	\argm{e}{E}{end frame number}{EOF}
	\argm{n}{N}{number of frame}{EOF}
	\argm{l}{L}{frame length.
                    Actually $\frac{L}{2}$ data are plotted.}{256}
	\argm{c, c2, c3}{\mbox{\em comment} 1 \sim 3}%
                        {comment for the graph}{N/A}
	\desc[1ex]{Usually, the options below do not need to be assigned.}
	\argm{W}{W}{width of the graph ($\times 100$ mm)}{0.25}
	\argm{H}{H}{height of the graph ($\times 100$ mm)}{1.5}
	\argm{z}{Z}{This option is used when data is written
                    recursively in the $y$ axis. the distance between
                    two graphs in the $y$ axis are given by $Z$. \\
                    If Z is not given, Z is as same as F}{}
	\argm{o}{xo \; yo}{origin of the graph.
                      if -o option exists, -O is not effective.}{95 30}
	\argm{g}{G}{type of frame of the graph ($0 \sim 2$) (see also \hyperlink{fig}{fig})}{2}
	\argm{cy}{cy}{first comment position}{-8}
	\argm{cy2}{cy2}{second comment position}{-14}
	\argm{cy3}{cy3}{third comment position}{-20}
	\argm{cs}{cs}{font size of the comments}{1}
	\argm{f}{f}{additional data file for \hyperlink{fig}{fig}}{NULL}
\end{options}

\begin{qsection}{EXAMPLE}
In this example, the magnitude of log spectrum is evaluated from
data in {\em data.f} file in float format, and the graph
with the running spectrum is sent in Postscript format to {\em data.ps}
file:
\begin{quote}
 \verb!frame < data.f | window |\! \\
 \verb!uels -m 15 | c2sp -m 15 |\! \\
 \verb!grlogsp | psgr > data.ps!
 \end{quote}
\end{qsection}

\begin{qsection}{SEE ALSO}
\hyperlink{fig}{fig},
\hyperlink{fdrw}{fdrw},
\hyperlink{xgr}{xgr},
\hyperlink{psgr}{psgr},
\hyperlink{glogsp}{glogsp},
\hyperlink{gwave}{gwave}
\end{qsection}
